\documentclass{article}
\usepackage{natbib}
\usepackage{url}
\begin{document}
\section{Introduction}
Transformers rely entirely on attention mechanisms and dispense with recurrence \citep{vaswani2017attention}.
Residual connections make it possible to train networks with more than a hundred layers \citep{he2016resnet}.
We train all models with the Adam optimizer \citep{kingma2014adam}.
Bidirectional pre-training of deep transformers improves many language understanding tasks \citep{devlin2019bert}.
Gradient folding reduces activation memory by 40\% during training \citep{chen2023gradfold}.
Randomly dropping units during training prevents co-adaptation of feature detectors and reduces overfitting \citep{smith2014dropout}.
The transformer architecture has since become the default for sequence modelling \citep{vaswani2017transformer, devlin2019bert}.
Fully convolutional networks produce dense pixel-wise predictions for semantic segmentation \citep{hochreiter1997lstm}.
Our implementation uses PyTorch \citep{pytorch}.
A standard textbook treatment of these methods is given by \citet{goodfellow2016deep}.
Large language models can perform tasks zero-shot without task-specific training \citep{radford2019gpt2}.
\section{Related work}
Normalizing layer inputs over each mini-batch speeds up the training of deep networks \citep[see][p.~3]{ioffe:2015-bn}.
Deep Q-networks reached human-level performance on many Atari 2600 games \cite{mnih2015dqn}.
\citet{maldacena1997} conjectured a duality between gravity in anti-de Sitter space and a conformal field theory on its boundary.
Kernel methods and support vector machines are treated in depth by \citet{scholkopf2002kernels}.
Hessian-free optimization and careful momentum schedules make recurrent networks trainable \citep{sutskever2013thesis}.
Generalization error falls as a power law in training set size \citep{hestness2017broken}.
Factorized convolutions and label smoothing improve Inception networks \citep{nocomma2016}.

\bibliographystyle{plainnat}
\bibliography{refs}
\end{document}
